\begin{tikzpicture}[
  >=Latex,
  line cap=round,
  line join=round,
  font=\small,
  eje/.style={->, black!75, line width=1.05pt},
  radial/.style={->, black!80, line width=1pt},
  centrifuga/.style={->, orange!85!black, line width=1.5pt},
  malla/.style={blue!45!black, line width=0.7pt},
  contorno/.style={blue!65!black, line width=1.15pt},
  formula/.style={draw=gray!55, rounded corners, fill=white,
    align=left, inner sep=8pt}
]
% Cilindro vertical en proyección oblicua: las elipses representan
% secciones circulares perpendiculares al eje de rotación.
\fill[blue!8] (-1.55,0.45) rectangle (1.55,4.35);
\fill[blue!10] (0,0.45) ellipse (1.55 and 0.48);
\fill[blue!16] (0,4.35) ellipse (1.55 and 0.48);
\draw[contorno] (-1.55,0.45) -- (-1.55,4.35);
\draw[contorno] (1.55,0.45) -- (1.55,4.35);
\draw[contorno] (0,4.35) ellipse (1.55 and 0.48);
\draw[contorno] (0,0.45) ellipse (1.55 and 0.48);

% Sección transversal donde se representan r y la aceleración centrífuga.
\fill[blue!22, opacity=0.72] (0,2.4) ellipse (1.55 and 0.48);
\draw[malla] (0,2.4) ellipse (1.55 and 0.48);
\draw[malla,densely dashed] (-1.55,2.4) -- (1.55,2.4);
\fill[black] (0,2.4) circle (2pt) node[below left=2pt] {$O$};
\fill[blue!65!black] (0.93,2.69) circle (2.2pt)
  node[above left=1pt] {$P$};
\draw[radial] (0.05,2.42) -- (0.90,2.68)
  node[midway,above left=1pt,fill=white,inner sep=1pt] {$\vec r=r\hat e_r$};

% Vectores centrífugos en la misma dirección radial; sus módulos crecen con r.
\draw[centrifuga] (0.34,2.51) -- (0.57,2.58);
\draw[centrifuga] (0.64,2.60) -- (1.00,2.71);
\draw[centrifuga] (0.93,2.69) -- (1.48,2.86)
  node[midway,above=2pt,fill=white,inner sep=1pt] {$\vec f(r)$};

% Eje de giro y vector de velocidad angular.
\draw[black!55,densely dashed] (0,0.15) -- (0,4.85);
\draw[eje] (-2.05,1.55) -- (-2.05,3.25)
  node[midway,left=3pt,fill=white,inner sep=1pt] {$\vec\Omega$};
\node[gray!75!black,align=center] at (0,4.98)
  {eje de rotación $z$};
\node[gray!75!black] at (0,-0.18) {sección transversal};

% Relación local entre aceleración ficticia y fuerza por unidad de volumen.
\node[formula,anchor=west] at (2.15,3.72) {
  En el marco que gira con el cilindro:\\[2pt]
  $\vec a_{\rm cf}=\Omega^2r\,\hat e_r$\\[4pt]
  $dm=\rho_0\,dV$\\[4pt]
  $\displaystyle \vec f=\frac{d\vec F_{\rm cf}}{dV}
    =\rho_0\vec a_{\rm cf}
    =\rho_0\Omega^2r\,\hat e_r$
};
\node[align=left,anchor=west,text width=5.0cm] at (2.15,1.75) {
  $\hat e_r$: dirección radial hacia fuera.\\
  En el eje, $r=0$ y $\vec f=\vec 0$;\\
  al aumentar $r$, aumenta $|\vec f|$.
};
\end{tikzpicture}